\documentclass{article}
\usepackage[T1]{fontenc}
\usepackage{lmodern}
\usepackage{graphicx}
\usepackage{amsmath,amssymb}
\usepackage[a4paper,margin=2cm]{geometry}
\usepackage[absolute,overlay]{textpos}
\setlength{\TPHorizModule}{1mm}
\setlength{\TPVertModule}{1mm}
\usepackage[indonesian]{babel}
\usepackage{hyperref}
\setlength{\emergencystretch}{2em}
% unit: Halaman 21

\begin{document}

% === Halaman 21 (python/native) ===

21

Kelemahan Mode ECB

1.

Karena plainteks sering mengandung bagian yang  berulang 

(sehingga terdapat blok-blok plainteks yang sama), maka 

enkripsinya menghasilkan blok cipherteks yang sama pula


➔ contoh berulang: spasi panjang

       ➔ mudah diserang secara statisitik dengan analisis frekuensi

\end{document}
